% teacher_preamble.tex — shared preamble for the chalk-teaching TEACHER'S NOTES.
% These are standalone, instructor-facing lesson plans (article class), separate
% from the student lecture notes.  Compile each teaching/lecNN_teacher_notes.tex
% directly with pdfLaTeX.
\usepackage[T1]{fontenc}
\usepackage[utf8]{inputenc}
\usepackage[margin=1in]{geometry}
\usepackage{amsmath,amssymb,amsthm,mathtools}
\usepackage{braket}
\usepackage{bm}
\usepackage{graphicx}
\usepackage{enumitem}
\usepackage{xcolor}
\usepackage{tikz}
\usetikzlibrary{arrows.meta}
\usepackage{quantikz}
\usepackage[most]{tcolorbox}
\usepackage[colorlinks=true,linkcolor=blue,urlcolor=blue]{hyperref}

% ---- Notation macros (consistent with the course preamble.tex) ----
\newcommand{\R}{\mathbb{R}}
\newcommand{\C}{\mathbb{C}}
\newcommand{\N}{\mathbb{N}}
\newcommand{\Z}{\mathbb{Z}}
\newcommand{\E}{\mathbb{E}}
\newcommand{\Prob}{\mathbb{P}}
\newcommand{\Hilb}{\mathcal{H}}
\newcommand{\ii}{\mathrm{i}}
\newcommand{\dd}{\mathrm{d}}
\newcommand{\ee}{\mathrm{e}}
\DeclareMathOperator{\tr}{tr}
\DeclareMathOperator{\Tr}{Tr}
\DeclareMathOperator{\Var}{Var}
\DeclareMathOperator{\diag}{diag}
\newcommand{\norm}[1]{\left\lVert #1 \right\rVert}
\newcommand{\abs}[1]{\left\lvert #1 \right\rvert}
\newcommand{\bigo}{\mathcal{O}}
\newcommand{\id}{\mathbb{I}}
\newcommand{\dg}{^{\dagger}}
% quantikz also defines \proj (a circuit element we do not use); override it with our projector.
\let\proj\relax
\newcommand{\proj}[1]{\ket{#1}\!\bra{#1}}
\newcommand{\Xgate}{\mathsf{X}}
\newcommand{\Ygate}{\mathsf{Y}}
\newcommand{\Zgate}{\mathsf{Z}}
\newcommand{\Hgate}{\mathsf{H}}
\newcommand{\CNOT}{\mathrm{CNOT}}
\newcommand{\Uf}{U_f}

% ---- Theorem-like environments (numbered per section, article style) ----
\theoremstyle{definition}
\newtheorem{definition}{Definition}[section]
\newtheorem{example}[definition]{Example}
\theoremstyle{plain}
\newtheorem{theorem}[definition]{Theorem}
\newtheorem{lemma}[definition]{Lemma}
\newtheorem{proposition}[definition]{Proposition}
\theoremstyle{remark}
\newtheorem{remark}[definition]{Remark}

% ---- Teaching-specific structure ----
% Session block header with a time budget, e.g. \block{Qubits and the Bloch sphere}{45}
\newcommand{\block}[2]{%
  \bigskip\noindent
  \begin{tcolorbox}[colback=blue!6,colframe=blue!55!black,boxrule=0.6pt,arc=2pt,
    left=8pt,right=8pt,top=5pt,bottom=5pt]
  \textbf{\large #1}\hfill\textbf{[\,#2\,min]}
  \end{tcolorbox}}

% Board-plan cue: what to write/draw on the board.
\newtcolorbox{board}{colback=gray!8,colframe=gray!55!black,title={On the board},
  fonttitle=\bfseries\small,boxrule=0.5pt,arc=2pt,left=6pt,right=6pt,top=3pt,bottom=3pt}

% Graphic-demonstration cue (what to sketch).
\newcommand{\sketch}[1]{\par\smallskip\noindent{\color{teal}\textbf{Sketch:}} #1\par\smallskip}

% Optional / advanced material that may be skipped (but indicated).
\newtcolorbox{advanced}{colback=orange!7,colframe=orange!65!black,
  title={Advanced --- may be skipped (mention only)},
  fonttitle=\bfseries\small,boxrule=0.5pt,arc=2pt,left=6pt,right=6pt,top=3pt,bottom=3pt}

% Timing hint inline, e.g. \tmark{10} for a ~10-minute chunk.
\newcommand{\tmark}[1]{{\small\color{blue!55!black}(\,$\approx$#1\,min)}}

% In-class problem (numbered per section) and a teacher-facing answer.
\newtheorem{inclass}[definition]{In-class problem}
\newenvironment{answer}{\par\smallskip\noindent\textit{Answer.}\ }{\par\smallskip}
